\section{Seam Winding as Source of the Integer Conditions}
\label{sec:seam-winding}

The integer triple $\bigl(3q_-+2q_+=0,\ q_-<0<q_+,\ \gcd(|q_-|,q_+)=1\bigr)$ entering input
(iii) of \Cref{thm:carrier-rigidity} is not assumed: it is the algebraic shadow of TFPT's seam
class $[u_\Sigma]\in H^1(\Sigma,\mathbb Z)$. The TFPT-side statement is the three-way obligation
\[
[u_\Sigma]=1
\;\Longrightarrow\;
\begin{cases}
3\,q_-+2\,q_+ = 0 & \text{(determinant trace-zero)},\\
q_- < 0 < q_+ & \text{(chirality orientation)},\\
\gcd(|q_-|, q_+) = 1 & \text{(primitivity of the seam class)}.
\end{cases}
\]
None of the three conjuncts is presented as a Lean theorem here; each is a Status-B obligation
of TFPT's seam-class theory. Their joint Lean target is
\[
\texttt{TFPT.Carrier.SeamWindingInterface.SeamWindingData}\ 3\ 2,
\]
which is the abstract typed bundle of the three integer conditions. The Lean-proved bridge
\[
\texttt{SeamWindingData}\ 3\ 2
\;\;\xrightarrow{\texttt{toPrimitiveOrientedDeterminantCarrier}}\;\;
\texttt{PrimitiveOrientedDeterminantCarrier}\ 3\ 2
\]
feeds the integer Diophantine input (iii) of \Cref{thm:carrier-rigidity}. The converse
\texttt{ofPrimitiveOrientedDeterminantCarrier} shows that the seam-winding interface is exactly
the existing rigidity input, renamed so that the upstream TFPT theorem
$[u_\Sigma]=1 \Rightarrow \texttt{SeamWindingData}$ has a Lean-typed conclusion.

\noindent
\textbf{What this section does not claim.}\ The topological identification
$[u_\Sigma]=1 \Rightarrow \texttt{SeamWindingData}\ 3\ 2$ remains a Status-B obligation
(open TFPT theorem). The present paper handles only the algebraic consequences once
\texttt{SeamWindingData} is supplied.

\section{Hypercharge and Gauge Quotient (Status C outlook)}
\label{sec:gauge-outlook}

\begin{tfptwarningbox}
The remainder of this paper presents the \emph{Standard-Model
packet outlook} conditional on \Cref{thm:carrier-rigidity}. These
sections are Status C (mathematical proposal): they are
\emph{not} part of the Lean-verified core and are
\emph{not} signature-locked in \texttt{AuditContract.lean}.
The text below indicates this consistently.
\end{tfptwarningbox}

The carrier normal form gives
\[
Y=\operatorname{diag}\left(-\frac13,-\frac13,-\frac13,\frac12,\frac12\right),
\]
and the internal stabilizer is conjectured to give the physical quotient
\[
G_{\mathrm{phys}}
=
\frac{SU(3)\times SU(2)\times U(1)_Y}{\mathbb Z_6}
\quad\text{\emph{(Status C: outlook, no Lean theorem)}}.
\]
The full proof obligation is to write a stabilizer theorem of the carrier datum, not to
reconstruct a known gauge group. The Lean-side typed target for that future theorem is expected
to read schematically
\[
\texttt{CarrierData}
\;\Rightarrow\;
\texttt{Z6QuotientStructure}\,Y,
\]
with \texttt{Z6QuotientStructure} the algebraic packaging of the stabilizer kernel; this
structure is \emph{not yet} introduced in the Lean note.

\section{One-Family Packet (Status C outlook)}
\label{sec:packet-outlook}

The spinor packet is conjectured to be
\[
S^+=\Lambda^{\mathrm{even}}E,
\qquad
\dim S^+=16
\quad\text{\emph{(Status C: outlook, no Lean theorem)}}.
\]
This section sketches the one-family decomposition including the right-handed neutrino state.
Character identities, trace identities, and the abelian index belong here. The expected Lean
target is a \texttt{DerivedSMRepresentations} module bundling
\texttt{LinearMap.dualHypercharge}, \texttt{ExteriorPower.hypercharge}, and
\texttt{TensorProduct.hypercharge}; not yet present in the Lean note.

\section{Discharging the Rank and Exterior Premises}

The editorial risk is that
\[
\dim \Lambda^2E_+=1,
\qquad
\Lambda^2E_-\cong E_-^\vee
\]
look like disguised Standard-Model inputs. In v10/v11 of the Lean note this risk is mitigated:
$\dim \Lambda^2 E_+ = 1$ follows from $\dim E_+ = 2$ via the Lean theorem
\texttt{exteriorPower.finrank\_eq}; and $\Lambda^2 E_- \cong E_-^\vee$ is the duality content of
Yukawa Stage~C (\texttt{PrimitiveIndecomposableYukawaCoupling}), itself derived from Stage~D under
$\texttt{Nontrivial}\,E_-$. The value-level transport Yukawa matrices belong downstream; the
upstream step here uses only the existence and operational predicates of $\omega$.

\section{Counting Outputs (Status C outlook)}
\label{sec:counting-outlook}

This paper outlines, conditional on the Status-A carrier core and on as-yet-open Status-B
content (representation theory of the carrier datum), the discrete counting outputs:
\[
N_{\mathrm{fam}}=3,
\qquad
\Oadm=48,
\qquad
N_\Phi=1,
\qquad
\bone=\frac{41}{10}.
\]
These are \emph{not} Lean theorems and \emph{not} precision-observable readouts. Their TFPT-side
derivation from carrier, seam winding, occupancy, and compact Higgs index closure is the
mathematical proposal of TFPT Paper 2's downstream layer.

\section{Main Result --- Carrier Rigidity (Status A)}

\Cref{thm:carrier-rigidity} is the main result of this paper. We restate it as a single
declarative summary:

\begin{quote}
\emph{Given the four explicit interfaces} --- $\epscar^2 = \mathbf 1$,
$E_+ \simeq K[X_0, X_1]_1$, an $\omega$ with $\textrm{CI} \wedge \textrm{CS}$ and
$\texttt{Nontrivial}\,E_-$, and the primitive oriented determinant-preserving carrier condition
on $(m, n) = (3, 2)$ --- \emph{Lean 4 proves} $(q_-, q_+) = (-2, 3)$, $\tr_E Y = 0$,
$6 Y^2 - Y - \mathbf 1 = \mathbf 0$, and
$Y = \operatorname{diag}(-\tfrac13,-\tfrac13,-\tfrac13,\tfrac12,\tfrac12)$.
\end{quote}

This is the entire load-bearing Status-A content of Paper 2. The Standard-Model packet
($S^+$, $G_{\mathrm{phys}}$, $N_{\mathrm{fam}}$, $\Omega_{\mathrm{adm}}$, $N_\Phi$, $b_1$) is
Status-C outlook conditional on this theorem and on as-yet-open Status-B obligations.

\section{Audit Protocol}

\noindent\textbf{Status A audit (Lean).}\ The Lean-verified core is audited by the seven-check
\texttt{audit.sh} gate of the companion Lean note. The audit gate covers: \texttt{lake build},
no \texttt{sorry}/\texttt{admit}, no \texttt{axiom}/\texttt{constant} declarations, no
\texttt{unsafe}/\texttt{opaque}/\texttt{partial}, every headline depends only on
$\{\texttt{propext},\,\texttt{Classical.choice},\,\texttt{Quot.sound}\}$, every headline
elaborates via \texttt{\#check} (\texttt{AuditCheck.lean}), and every headline is signature-locked
via $\texttt{example : <exact intended type>}$ (\texttt{AuditContract.lean}). Current state:
\texttt{AUDIT:~PASS}.

\smallskip
\noindent\textbf{Editorial audit (this paper).}

\begin{center}
\footnotesize
\begin{tabularx}{\textwidth}{@{}p{0.18\textwidth}p{0.08\textwidth}Y@{}}
\toprule
\textbf{Audit item} & \textbf{Tier} & \textbf{Question} \\
\midrule
Rank discharge
  & A
  & Do Higgs Stage B and Yukawa Stage D force $(\dim E_-,\dim E_+)=(3,2)$ \emph{before} the carrier polynomial is written? \\
Carrier polynomial
  & A
  & Is $6Y^2-Y-\mathbf 1=0$ presented only as the minimal polynomial of the derived roots
    $\{-1/3, 1/2\}$? \\
Hypercharge
  & A
  & Is the determinant normalisation fixed before any phenomenological matching? \\
Seam-winding interface
  & B
  & Is the integer triple $(\textrm{trace-zero},\textrm{orient.},\textrm{prim.})$ presented as the
    Lean-typed conclusion of an open TFPT theorem
    $[u_\Sigma]=1 \Rightarrow \texttt{SeamWindingData}$? \\
Yukawa interface
  & A/B
  & Is the rank-3 discharge presented as a Lean theorem about $\omega$ with operational
    predicates, with the canonicity of $\omega$ from TFPT data marked as a Status-B obligation? \\
Packet ($S^+$)
  & C
  & Is $S^+=\Lambda^{\mathrm{even}}E$ marked as Status-C outlook, not as a Lean theorem? \\
Gauge quotient
  & C
  & Is the $\mathbb Z_6$ quotient marked as a Status-C outlook with an explicit Lean-side
    target name (\texttt{Z6QuotientStructure}) for the future theorem? \\
Counting outputs
  & C
  & Are $N_{\mathrm{fam}},\Omega_{\mathrm{adm}},N_\Phi,b_1$ presented as Status-C readouts, not
    as Lean theorems? \\
\bottomrule
\end{tabularx}
\end{center}

\clearpage
\section{Main Technical Development}

This section contains the main technical development assigned to this paper by the TFPT 4.5 clean split. Cross-paper background is referenced through dependency and scope boxes; extended backend material is kept in the Technical Companion.

\begin{tfptwarningbox}
\textbf{Status-tier mapping for the source-extract body.}\
The text below is the source-manuscript extract for the Paper-2
clean split. It uses the legacy \LaTeX{} \texttt{theorem}
environment uniformly across content of \emph{different} formal
status in the new five-tier taxonomy. The body is to be read
under the following mapping:

\smallskip
\begin{itemize}[leftmargin=1.7em, itemsep=0pt]
    \item Carrier-core algebraic statements
        (carrier polynomial, $\tr_E Y = 0$, $(q_-,q_+)=(-2,3)$,
        carrier involution to projectors, integer rigidity)
        carry status \stag{L} via the Lean theorems referenced
        in \Cref{thm:carrier-rigidity}.
    \item Internal-stabilizer / $\mathbb Z_6$ quotient theorems
        (e.g.\ "Internal carrier stabilizer",
        "Faithful physical global gauge group"),
        one-family packet / $S^+ = \Lambda^{\mathrm{even}} E$
        theorems, the joint-discrete-admissible-sector theorem,
        the intrinsic-branch / canonical-rigid-object
        statements, and the manuscript-level "carrier rigidity
        from the primitive component of the branch Yukawa
        tensor" carry status \stag{M} or \stag{C}:
        \begin{itemize}[leftmargin=1.2em, itemsep=0pt]
            \item \stag{M} where the source manuscript provides
                a complete proof at standard mathematical-writing
                rigour (e.g.\ stabilizer / $\mathbb Z_6$
                construction, exterior packet calculation, joint
                discrete admissible sector).
            \item \stag{C} where the statement is conditional on
                upstream open content (representation theory of
                derived powers, anomaly bookkeeping that depends
                on \stag{C} occupancy values, etc.).
        \end{itemize}
    \item Strong-CP, hadronic admissibility, and seam-inflow
        statements that depend on the analytic hypotheses of
        Paper 4 carry status \stag{C} (and are imported by
        Paper 4 with explicit conditional framing).
    \item The "Yukawa generator" / "primitive Yukawa tensor"
        wording in the source extract is the manuscript
        precursor of the v11 Lean Stage~D interface
        $(\omega, \texttt{CI}, \texttt{CS}, \texttt{Nontrivial})$.
        The bridge
        \begin{quote}\small
        branch Yukawa tensor (manuscript)
        \;$\Rightarrow$\;
        \texttt{BoundaryYukawaKernel}
        \;$\Rightarrow$\;
        $\omega$ with CI, CS, Nontrivial \stag{I}
        \end{quote}
        is intended; the body retains the manuscript wording.
\end{itemize}

\smallskip
None of the body theorems below is \stag{L} unless explicitly
named in \Cref{thm:carrier-rigidity}. Section
\S\ref{sec:formal-status} of this paper is the binding
statement of \stag{L} content; the body extract is the source
manuscript for the \stag{M} / \stag{C} content listed in the
\stag{C} outlook sections
(\Cref{sec:gauge-outlook,sec:packet-outlook,sec:counting-outlook}).
\end{tfptwarningbox}

\tfptDraftXrefNotice

\input{source_extracts/02_carrier_source.tex}

\section{Source Extraction Map}

\begin{tfptsourcebox}
Use \sourceDraft:
\begin{itemize}[leftmargin=1.5em]
    \item Sections 3.3, 3.4, and 3.5 for carrier normal form, internal stabilizer, one-family
    packet, character identities, trace identities, and abelian index.
    \item Sections 5.1--5.3 only for carrier-relevant compression identities.
    \item Sections 6.1--6.6 for families, occupancy, and compact Higgs index.
    \item Sections 8.11--8.15 for the Yukawa-based discharge of carrier premises.
    \item Section 8.19 and Section 9.11 only for the joint discrete admissible sector.
\end{itemize}
\end{tfptsourcebox}

\begin{tfptwarningbox}
The paper stands or falls on \Cref{thm:carrier-rigidity} (Status A, Lean-verified) and on the
honest separation of Status-A core from Status-C outlook. Exact electromagnetic numerics,
CKM/PMNS, masses, OS reconstruction, cosmology, and E8 are deferred to later papers (and remain
Status C until separately formalised). The Status-C sections (\Cref{sec:gauge-outlook},
\Cref{sec:packet-outlook}, \Cref{sec:counting-outlook}) must be read as
\emph{conditional} on the Status-A core, not as additional Lean theorems.
\end{tfptwarningbox}


\begin{tfptexports}
\textbf{Status A (Lean theorems).}\ $(q_-, q_+) = (-2, 3)$;\ \ $\tr_E Y = 0$;\ \
$6 Y^2 - Y - \mathbf 1 = \mathbf 0$;\ \
$Y = \operatorname{diag}(-\tfrac13,-\tfrac13,-\tfrac13,\tfrac12,\tfrac12)$;\ \ the carrier datum
$E = E_- \oplus E_+$ with $(\dim E_-, \dim E_+) = (3, 2)$.

\smallskip
\textbf{Status B (Lean-typed targets for open TFPT theorems).}\
\texttt{BoundaryPolarization} (analytic Paper-1 content);\ \
\texttt{BoundaryYukawaKernel} (TFPT-canonical $\omega$);\ \
\texttt{SeamWindingData} (seam-class topological content);\ \
\texttt{HiggsSchemeCohomologyShadow} (already Lean-proved $\to$ \texttt{DegreeOneTwoVarForm};
remaining Proj-side identification is general algebraic geometry).

\smallskip
\textbf{Status C (mathematical proposal, downstream outlook).}\
$S^+ = \Lambda^{\mathrm{even}} E$;\ \
$G_{\mathrm{phys}} = \big(SU(3)\times SU(2)\times U(1)_Y\big) / \mathbb Z_6$;\ \
$N_{\mathrm{fam}} = 3$;\ \ $\Oadm = 48$;\ \ $N_\Phi = 1$;\ \ $b_1 = 41/10$.
\end{tfptexports}

\section{Not Used Here}

Exact $\alpha$ numerics, CKM/PMNS closure, pole-mass ledgers, OS/CAR reconstruction, local
Minkowski nets, scattering, gravity/metrology, cosmology, CMB targets, and E8 grammar are not used
as proof inputs in this paper.

\end{document}
